\section{Introduction}
\label{sec:intro}

When Ward Cunningham likened expedient code to borrowed money, he gave software
engineering a durable insight: some costs are not avoided but \emph{deferred}, and
deferral has a price~\cite{cunningham1992}. Three decades of research have developed
that insight into a family of debt constructs---code, design, architectural,
documentation, and security debt among them~\cite{kruchten2012,li2015,tom2013,
alves2016,rindell2019}---each naming a way in which present convenience becomes
future engineering cost.

This paper investigates a deferred cost that the debt family does not yet name, and
that cloud-native operations manufacture at scale. When an infrastructure change is
made---a deployment, a configuration edit, a permission grant, an emergency fix---an
organization either captures durable evidence of the change's essential facts (who
acted; what changed; why; under which authorization; with what outcome; how the
resulting state can be reconstructed) or it does not. Not capturing them is almost
always the locally rational choice: capture costs effort now, while the missing
record costs nothing today. The cost arrives later, and in a different budget: during
the incident whose responders must establish what changed last Tuesday; during the
audit that must connect running workloads to authorized changes; during the security
investigation that must distinguish an operator's forgotten hotfix from an intrusion;
after the departure of the engineer who was the only index into a year of undocumented
decisions; during disaster recovery, when ``how the system can be rebuilt'' turns out
to live nowhere but in the system being rebuilt; and in litigation or regulatory
inquiry, where the absence of records is itself a finding~\cite{sedona2018,sox2002}.

We name the accumulated form of this deferred cost \emph{evidence debt}, and we spend
the paper making the name earn its keep. The financial metaphor is seductive and
mostly wrong in its details, and one of our contributions is to say precisely where
(\cref{sec:model-anti,sec:disc-metaphor}). What survives scrutiny is a phenomenon
with its own structure: evidence gaps are created by identifiable practices at change
time (\cref{sec:conditions}); they are syntactically observable as \emph{deficits}
against a declared evidence schema (\cref{sec:definitions}); their economic meaning
is the \emph{excess reconstruction burden}---effort, accepted wrong answers, and
policy abstention---they impose on future queries against
operational history (\cref{sec:model}); that cost grows over time through
identifiable decay processes---log expiration, ticket deletion, personnel turnover,
resource ephemerality, identifier loss---whose superposition produces plateau-and-cliff
effort and abstention growth under sound exhaustive policies rather than compounding
(\cref{sec:interest}); and past a knowable boundary the history
becomes not expensive but \emph{irrecoverable}, a mandatory write-off with no
financial analogue (\cref{sec:model-anti}).

The central claim is treated as a falsifiable hypothesis rather than as an
assumption: \emph{organizations can defer evidence capture, but
they cannot eliminate the eventual cost of reconstructing missing operational
history}. As stated, the claim is not even obviously true---redundancy across
operational data sources might, in principle, make most deferred capture free,
because the history could be rebuilt from what other systems recorded anyway. Our
simulation study (\cref{sec:simulation}) shows that this objection has real force at
low degradation rates (redundancy absorbs single-class link loss almost completely)
and fails structurally at higher rates and under record deletion, where
reconstruction collapses, and---more insidiously---that a policy accepting unique
heuristic candidates still returns wrong history at rates reaching 30.3\% of
accepted links under dense record deletion. Acceptance is decided from the corpus
\emph{before} ground truth is consulted, so lucky rejected guesses cannot become
supported answers after the fact. Deferred capture is thus repaid in one of three
policy-burden currencies: effort, accepted error, or abstention; certified physical
irrecoverability is measured separately.

\subsection{Research Questions}
\label{sec:rqs}

\begin{description}[leftmargin=3.2em,itemsep=1pt,topsep=2pt]
\item[RQ:] Can accumulated deficiencies in operational evidence be formally modeled
and measured as evidence debt?
\item[RQ1:] What conditions create evidence debt?
\item[RQ2:] How can evidence debt be quantified?
\item[RQ3:] How do evidence deficits alter the effort, accuracy, and confidence of
later reconstruction?
\item[RQ4:] Under which conditions does evidence debt become irrecoverable rather
than merely expensive?
\end{description}

RQ1 is answered analytically by the debt-creating-conditions taxonomy of
\cref{sec:conditions}; RQ2 by the definitions, model, and metric suite of
\cref{sec:definitions,sec:model,sec:metrics}; RQ3 and RQ4 jointly by the model's
predictions (\cref{sec:model,sec:interest}) and their test in the executed
simulation study (\cref{sec:simulation}), with external validity addressed by the
bounded repository observation and full field protocol of \cref{sec:field}.

\subsection{Contributions and Epistemic Status}
\label{sec:contributions}

\begin{enumerate}[label=\textbf{C\arabic*.},leftmargin=2.2em,itemsep=1.5pt,topsep=2pt]
\item \textbf{Concept and definitions} (conceptual): non-circular formal definitions
of evidence deficit, evidence debt, principal, interest, liability, reconstruction
cost, and irrecoverability, each anchored to an observable or measurable property and
none defined in terms of the financial metaphor (\cref{sec:definitions}).
\item \textbf{Formal model} (conceptual/analytic): evidence debt as expected excess
effort, accepted-error, and abstention burden over a query workload, with certified
irrecoverability as a distinct evidence state; reconstruction as record-linkage
inference over a fragmented evidence corpus~\cite{fellegi1969,christen2012}; decay
as per-source survival processes (\cref{sec:model}).
\item \textbf{Anti-metaphor results} (analytic): propositions showing that
reconstruction effort and abstention exposure under sound exhaustive policies are stochastic and plateau-and-cliff
rather than compound, that
the debt is bounded by mandatory write-offs, and that the constructs remain
well-defined with the metaphor deleted (\cref{sec:model-anti,sec:disc-metaphor}).
\item \textbf{Metric suite with rejection} (methodological): six candidate metrics
critically evaluated; one (a scalar evidence-debt interest rate) is rejected as
ill-defined under the model and replaced by hazard-based alternatives
(\cref{sec:metrics}).
\item \textbf{Executed simulation study} (empirical, synthetic): a reproducible
evidence-degradation experiment---complete synthetic event chains; five evidence
classes degraded at 1--40\%; a pairwise coverage arm plus a paired link-level path
census that tests the formal interaction premise directly; removal and sensitivity arms
(redundancy off, heuristic-window stress, actor-assumption breakage, single-knob
density sweep) that separate design-entailed demonstrations from falsifiable
tests; twenty seeds; all code and data included in the accompanying artifact, with results tables generated
programmatically from run outputs---including an end-to-end computation of
$\ED_\rho(t)$ under a declared workload (\cref{sec:simulation}). The study
validates that the constructs are computable and discriminating; it estimates no
real-world magnitudes.
\item \textbf{Public-repository observation} (empirical, bounded): a reproducible
age-stratified sample of 200 merged Argo CD pull requests and 30 releases under a
narrow, machine-testable repository schema, with authorship/change-type strata and
a seeded 50-PR manual audit. It establishes that missing external issue linkage
occurs outside the simulation, while review-record and release-lineage presence
controls remain discriminating; it does not measure intent adequacy, decay, debt,
or cost (\cref{sec:field}).
\item \textbf{Field protocol} (methodological, not executed): instruments for a
deficit census, reconstruction drills, and longitudinal decay observation in real
organizations (\cref{sec:field}). The present article deliberately establishes the
theoretical and synthetic foundation; it does not claim field calibration of ERC,
FRR, AR, or consequence weights.
\end{enumerate}

Throughout, we separate established results (cited), analytic claims (proved within
the model), empirical claims (measured and scoped to their source), hypotheses
(labeled), and proposals (labeled). The only empirical statement about a real
project is the bounded public-repository observation of \cref{sec:field}; it is
explicitly separated from evidence-debt and cost measurement.

\subsection{Scope and Non-Goals}

We study operational evidence for infrastructure and configuration change in
cloud-native systems. We do not propose a logging product, do not claim that more
retention is always better (retention has cost and legal risk~\cite{gdpr,nist80092}),
and do not address the orthogonal question of whether captured evidence is
trustworthy against tampering---the secure-logging literature covers
integrity~\cite{schneier1999,crosby2009,haber1991}; our subject is \emph{existence,
connectivity, and survivability} of evidence, which integrity mechanisms presuppose.
